\section{Background and notation}\label{sec:background}

\paragraph{Schemes.}
Index the entries of an $n\times m$ matrix $A$ by pairs $(i,j)$ and identify the pair with the coordinate $i\,m+j$
(row-major). Let $e_{(i,j)}$ be the corresponding standard basis vector of $\F^{nm}$. The matrix-multiplication tensor is
\[
  T_{n,m,p} \;=\; \sum_{i<n,\;j<m,\;k<p} e_{(i,j)}\tensor e_{(j,k)}\tensor e_{(i,k)}
  \;\in\; \F^{nm}\tensor\F^{mp}\tensor\F^{np}.
\]
Note that the third factor is indexed by the entry $(i,k)$ of the \emph{product}, not transposed.
A \emph{scheme of rank $r$ for $\MM nmp$ over $\F$} is a list of $r$ triples $(u_t,v_t,w_t)$ of nonzero vectors with
$\sum_t u_t\tensor v_t\tensor w_t = T_{n,m,p}$. Equivalently, for all $A\in\F^{n\times m}$ and $B\in\F^{m\times p}$,
\[
  (AB)_{ik} \;=\; \sum_{t=1}^{r} w_t[(i,k)]\;\langle u_t,A\rangle\,\langle v_t,B\rangle ,
\]
so the scheme computes the product with $r$ multiplications of linear forms in the entries of $A$ and $B$.
We write $\rk_{\F}\MM nmp$ for the least such $r$. Every scheme found by any of our tools is an \emph{upper bound}
for this quantity; we never certify a lower bound.

\paragraph{Orientation.}
The tensor $T_{n,m,p}$ is carried to $T_{m,p,n}$ and to $T_{p,m,n}$ by cyclic shifts of the three factors and by transposition of the
matrices, with the appropriate re-indexing of coordinates. Hence $\rk_{\F}\MM nmp$ depends only on the multiset $\{n,m,p\}$. All tables and
all price computations in this paper key on the \emph{sorted} shape, and a scheme found for one orientation is turned into a scheme for
another by an explicit coordinate permutation (with a transposition of a factor when the permuted edge is reversed) that is itself
covered by the exactness gate of \cref{sec:repr}.

\paragraph{Symmetries.}
The isotropy group $\mathrm{GL}_n(\F)\times\mathrm{GL}_m(\F)\times\mathrm{GL}_p(\F)$ acts on schemes for $\MM nmp$ and is the source of
the many equivalent presentations of a fixed scheme \cite{degroote1978isotropy,burichenko2015symmetries}.
The walker of \cref{sec:walker} does not use it; the derivation step of \cref{sec:derive} acts on presentations by a different mechanism
(refactoring the terms that share a factor).

\paragraph{Flips, reductions and plus transitions.}
Regard a scheme as a set of terms, with the convention that two equal terms cancel (characteristic two) and a term with a zero
factor is dropped. The following identity is the basis of the flip graph of Kauers and Moosbauer \cite{kauers2023flip}.

\begin{lemma}[flip]\label{lem:flip}
For vectors $s$, $l_1,l_2$, $r_1,r_2$ of the three factor spaces (in any assignment of the roles to the three factors),
\[
  s\tensor l_1\tensor r_1 + s\tensor l_2\tensor r_2
  \;=\; s\tensor l_1\tensor (r_1+r_2) + s\tensor (l_1+l_2)\tensor r_2
  \qquad\text{over }\F .
\]
\end{lemma}
\begin{proof}
The right-hand side expands to $s\tensor l_1\tensor r_1 + s\tensor l_1\tensor r_2 + s\tensor l_1\tensor r_2 + s\tensor l_2\tensor r_2$, and the
two middle terms cancel in characteristic two.
\end{proof}

Two terms that share a factor can therefore be replaced by two other terms without changing the sum. The replacement never
increases the number of terms; it lowers it when one of the new terms has a zero factor (which happens when the two old terms
already share a second factor: the pair \emph{merges}) or equals a third term of the scheme (both cancel). Such coincidences are called
\emph{reductions}; a flip walk moves through schemes of equal rank until one occurs. A rank-\emph{increasing} move is needed to
leave a basin: the \emph{plus transition} of adaptive flip graphs \cite{arai2024adaptive} raises the rank by one by splitting a factor, taking the new part $e$ from another term.
Our walker's rank-raising move is a variant: it splits one factor with a random unit vector $e$,
$f\tensor g\tensor h = e\tensor g\tensor h + (f+e)\tensor g\tensor h$, raising the rank by one.
Our walker uses flips and these splits and nothing else (\cref{sec:walker}).

\paragraph{Related searches.}
Flip graphs were introduced for formats up to $5\times5\times5$ with lifting from $\F$ to $\mathbb Z$ by Hensel's lemma \cite{kauers2023flip},
extended to formats $(n,m,6)$ \cite{kauers2025six} and to formats with all dimensions up to 8 by a meta flip graph
\cite{kauers2025meta}, restricted to symmetric schemes \cite{moosbauer2025flip}, run over ternary coefficients on GPUs
\cite{perminov2025ternary} and in a CPU framework up to $16\times16\times16$ \cite{perminov2026framework}, combined with the serendipitous product (Smith; generalized by Sedoglavic)
\cite{perminov2026serendipitous}, and run over $\F$ on a GPU to reach rank 245 for $7\times7\times7$ \cite{medley2026gpu}.
Reinforcement learning \cite{fawzi2022discovering}, SAT solvers \cite{heule2019local,heule2021new} and evolutionary program search
\cite{novikov2025alphaevolve} found other records.
Our pipeline differs in target and in emphasis: it works on formats up to $32\times32\times32$ from seeds that already exist, and the
part of it that is new to the literature (as far as we know) is not the moves (projection between formats is an edge of the meta flip graph \cite{kauers2025meta}, and the public catalogues \cite{sedoglavic2026catalogue,lacelle2026catalog} already derive ranks exhaustively by composition and projection) but the shared-factor basis rewrite applied before projection, the pricing and materialization of compositions from GF(2)-valid leaves with hash-pinned witnesses, the replay of public derivations over $\F$, and verifiable provenance.
